% =====================================================================
% draft_sections/scope.tex
% Honest "When does the framework help?" scope analysis.
% 2026-07-23 revision: the earlier traffic-speed boundary claim was an artifact
% of missing-value contamination (METR-LA zeros); on cleaned data the masks help
% on speed networks too (metrla pl96: 0.706 vs 0.915 all-ones, 3 seeds). The
% StructMamba all-ones ablation (16/16, +13.8%) replaces the "neutral" claim.
% All numbers recomputed from audit.json.
% =====================================================================

\subsection{When Does the Framework Help?}
\label{sec:scope}

A useful account of a mechanism states both axes of its scope: the \emph{data} on
which it helps, and the \emph{backbone} into which it is dropped. Our experiments
separate the two cleanly.
On the data axis, MRMA helps precisely when the training split admits three
\emph{complementary} relation graphs that recover a genuine local dependency
structure. This holds across large traffic sensor networks of both physical
quantities we tested: traffic \emph{flow} (all four PEMS benchmarks) and traffic
\emph{speed} (METR-LA, PEMS-BAY), where directed segment coupling, marginal
congestion co-movement, and lead-lag propagation zones are physically distinct
phenomena, so the partial-correlation, Pearson, and 2D-structural-entropy
community graphs each carry signal the others miss. On METR-LA at pred-96, for
example, the masked model reaches 0.706 MSE against 0.915 for the identical model
with all-ones masks ($+22.8\%$, 3 seeds)---provided missing readings (encoded as
zeros) are interpolated before both training and graph estimation. We flag this
as a practical caveat rather than a boundary: on the \emph{contaminated} series
the estimated graphs misalign and the same masks turn neutral-to-harmful, so
graph quality is a data-hygiene requirement of the method.
On the backbone axis, the masks deliver their full benefit wherever channel
mixing is otherwise unstructured---all 16 PEMS cells on \itrans{}'s single global
attention ($+9.7\%$) \emph{and} all 16 on a gated state-space recurrence
(\textsc{StructMamba}, $+13.8\%$)---but only partially (5 of 16 cells) inside
CrossFormer's two-stage router, which already imposes a learned channel
structure that the fixed masks compete with rather than replace.

Outside this regime the framework neither helps nor materially hurts---it
converges to the backbone---and we report the boundary explicitly rather than
suppressing it.
Two conditions mark the edge.
\emph{(i) Diffuse or single-driver dependency:} when one dominant factor drives
all channels (photovoltaic panels under shared irradiance; finance series such
as exchange-rate and NASDAQ equity data with pervasive factor co-movement), the
three graphs collapse toward the same structure, the relational heads lose their
division of labour, and the masked model tracks full attention within noise.
\emph{(ii) Low channel count:} for small-$N$ data (ETT-family, Weather, with
$N{\lesssim}20$) the top-$k$ neighbourhoods already cover most channel pairs, so
there is little to sparsify and the mask union is near-dense; gains fall within
seed noise.
This is the same mechanism read in reverse: the benefit is contingent on
graph--data alignment, so where alignment fails---too few channels, one global
driver, or corrupted estimation data---the benefit vanishes with it.

We regard this two-axis scope as a feature of an honest paper, not a caveat to be
minimised.
MRMA is a \emph{structural prior for large sensor networks whose host mixes
channels without its own learned structure}: it is exactly as strong as the
structure its three graphs can recover from the training split and as the room
the host backbone leaves for a static mask to fill---no stronger, and, on the
datasets or backbones where either condition fails, no weaker than the backbone
it drops into.
